\documentclass[11pt,a4paper]{report}

\usepackage[utf8]{inputenc}
\usepackage[T1]{fontenc}
\usepackage[spanish,es-nodecimaldot]{babel}
\usepackage{lmodern}
\usepackage{microtype}
\usepackage[a4paper,margin=2.5cm]{geometry}
\usepackage{xcolor}
\usepackage{booktabs}
\usepackage{longtable}
\usepackage{tabularx}
\usepackage{array}
\usepackage{graphicx}
\usepackage{float}
\usepackage{hyperref}

\hypersetup{
  colorlinks=true,
  linkcolor=blue!50!black,
  urlcolor=blue!50!black,
  pdftitle={Informe tecnico de implementaciones desde f81eb49},
  pdfauthor={Codex}
}

\renewcommand{\arraystretch}{1.18}
\emergencystretch=2em

\newcommand{\commit}[1]{\texttt{#1}}
\newcommand{\code}[1]{\nolinkurl{#1}}
\newcommand{\story}[1]{\texttt{#1}}
\newcommand{\confidence}[1]{\textbf{#1}}
\newcommand{\qafigure}[2]{
  \begin{figure}[H]
    \centering
    \includegraphics[width=0.92\textwidth,height=0.78\textheight,keepaspectratio]{#1}
    \caption{#2}
  \end{figure}
}

\title{
  \Huge Informe t\'ecnico de implementaciones\\[0.4em]
  \Large desde el commit \commit{f81eb49f7c6a90f983f67fd831cf13872ce5326c}
}
\author{Reconstrucci\'on sobre el repositorio \code{paw-2026a-01}}
\date{19 de abril de 2026}

\begin{document}

\maketitle
\clearpage
\pagenumbering{roman}
\tableofcontents
\clearpage
\pagenumbering{arabic}

\chapter{Introducci\'on y criterio de reconstrucci\'on}

\section{Objetivo}

Este documento resume y explica, con foco t\'ecnico, las implementaciones incorporadas en el repositorio \code{paw-2026a-01} a partir del commit base \commit{f81eb49f7c6a90f983f67fd831cf13872ce5326c} del 12 de abril de 2026. El objetivo no es listar commits en bruto, sino reconstruir qu\'e capacidades nuevas quedaron disponibles, c\'omo fueron implementadas y qu\'e capas del sistema se modificaron para soportarlas.

\section{Fuentes utilizadas}

La reconstrucci\'on se hizo combinando cinco fuentes:

\begin{itemize}
  \item historial Git entre \commit{f81eb49} y \commit{72a5d7d8};
  \item la wiki local en \code{PAW-Wiki/docs/}, en particular \code{index.md}, \code{log.md}, \code{wiki/resumen-notas-sprint-2.md}, \code{wiki/plan-implementacion-reservas.md} y \code{wiki/2026-04-17_cierre-implementacion-reservas_v1.md};
  \item las capturas del tablero de Plane del Sprint 2 provistas fuera del repositorio;
  \item los logs locales de runtime bajo \code{logs/};
  \item el c\'odigo actual en \code{models}, \code{persistence}, \code{services} y \code{webapp}.
\end{itemize}

\section{Delimitaci\'on del rango}

El rango analizado contiene:

\begin{itemize}
  \item 93 commits posteriores al baseline;
  \item 9 merges;
  \item 267 archivos tocados;
  \item 29991 inserciones y 3471 eliminaciones.
\end{itemize}

Eso obliga a organizar el informe por \'epicas funcionales y no por orden cronol\'ogico estricto.

\section{Advertencia de atribuci\'on}

No todas las tarjetas visibles en Plane equivalen a una implementaci\'on completa nacida dentro de este rango. En algunos casos, el trabajo del rango consiste en:

\begin{itemize}
  \item terminar una implementaci\'on que ya hab\'ia empezado antes;
  \item mover una capacidad a nuevas rutas o nuevas vistas;
  \item endurecer validaciones, seguridad o i18n;
  \item corregir regresiones sobre una feature preexistente.
\end{itemize}

Por eso, cuando una story no puede atribuirse con confianza alta al rango \commit{f81eb49..HEAD}, el informe lo marca expl\'icitamente como refinamiento o cierre parcial.

\chapter{Panorama general del trabajo realizado}

\section{Resumen ejecutivo}

Desde el commit base se consolidaron siete l\'ineas principales de trabajo:

\begin{enumerate}
  \item cierre del ciclo completo de autenticaci\'on y cuenta de usuario;
  \item ampliaci\'on del perfil reviewer y de las listas personales;
  \item dashboard owner y backoffice admin para moderar restaurantes;
  \item implementaci\'on integral del sistema de reservas;
  \item consolidaci\'on del mailing funcional y de los procesos background;
  \item mejoras de UX en explore, paginaci\'on y formularios;
  \item hardening transversal de seguridad, validaci\'on, arquitectura y tests.
\end{enumerate}

\section{Commits estructurantes}

\begin{longtable}{p{0.18\textwidth}p{0.74\textwidth}}
\toprule
\textbf{Commit} & \textbf{Aporte estructural} \\
\midrule
\endhead
\commit{f497af22} & alta de Spring Security, login, claim de cuenta y flujo base de auth. \\
\commit{327aeeb7} & reset de password para cuentas reclamadas. \\
\commit{03bf6161} & owner dashboard, moderaci\'on admin y reviews autenticadas. \\
\commit{37e7997e} & modelo, schema y DAOs base del sistema de reservas. \\
\commit{03fd92b8} & flujo de creaci\'on de reservas desde la p\'agina del restaurante. \\
\commit{5387190a} & gesti\'on owner de reservas con filtros, reasignaci\'on y acciones por estado. \\
\commit{155ab1b2} & scheduler para expirar pendientes y completar confirmadas. \\
\commit{3b6c5221} & mails transaccionales y recordatorio diario. \\
\commit{2e16c230} & reviews verificadas asociadas a reservas completadas. \\
\commit{4e884b77} & registro unificado en \code{/register}. \\
\commit{10915c6a} & verificaci\'on de cuenta por correo con autologin. \\
\commit{2e52cd5a} & separaci\'on de edici\'on de restaurante, men\'u y reservas. \\
\commit{28014edf} & wizard en la edici\'on/creaci\'on de restaurante. \\
\commit{5974e49e} & filtros y paginaci\'on para reservas y restaurantes guardados. \\
\commit{120e0432} & refactor para volver a sacar l\'ogica de negocio fuera de controllers. \\
\bottomrule
\end{longtable}

\chapter{Autenticaci\'on y ciclo de cuenta}

\section{Base del flujo de autenticaci\'on}

El commit \commit{f497af22} instal\'o la primera base consistente del dominio de autenticaci\'on. El trabajo se reparti\'o en tres capas:

\begin{itemize}
  \item en \code{webapp}, con \code{WebAuthConfig}, \code{PawAuthUser}, \code{PawUserDetailService}, \code{AccessHelper} y \code{AuthController};
  \item en formularios y validadores, con \code{UniqueEmail}, \code{UniqueUsername} y \code{PasswordsMatch};
  \item en vistas, con las pantallas de login, claim y registro owner/reviewer.
\end{itemize}

La idea principal fue dejar de modelar auth como una suma de pantallas aisladas y pasar a un flujo de seguridad centralizado en Spring Security, con principal tipado y helpers de acceso reutilizables.

\section{Mejoras de UX y consolidaci\'on visual}

Los commits \commit{be6ebcf6}, \commit{dfac52c9} y \commit{e335e81d} no agregan reglas nuevas de negocio, pero s\'i son importantes porque convierten el flujo de auth en una experiencia coherente:

\begin{itemize}
  \item introducen \code{authSplitLayout.tag};
  \item reorganizan \code{login.jsp}, \code{register.jsp}, \code{claim-activate.jsp}, \code{claim-request.jsp} y sus variantes;
  \item alinean la UI de login, claim y registro con un lenguaje visual com\'un;
  \item agregan activos gr\'aficos espec\'ificos para registro.
\end{itemize}

Esto importa porque en el Sprint 2 varias correcciones no eran de l\'ogica sino de claridad visual, densidad del formulario y navegaci\'on entre pantallas.

\section{Login por email}

El commit \commit{2766ccc2} expande el criterio de autenticaci\'on para permitir login por email adem\'as de username. La implementaci\'on se hizo dentro de \code{PawUserDetailService}, no duplicando controllers ni creando una segunda ruta de autenticaci\'on.

La decisi\'on de dise\~no fue correcta por dos motivos:

\begin{itemize}
  \item mantiene a Spring Security como due\~no del login;
  \item encapsula la resoluci\'on email/username en el mismo servicio que ya cargaba el principal.
\end{itemize}

\section{Claim de cuentas legacy y reset de password}

El rango tambi\'en cerr\'o el ciclo de cuentas ya existentes. Primero se consolid\'o el claim de reviewer legacy; luego, con \commit{327aeeb7}, se agreg\'o reset de password para cuentas reclamadas.

El recorrido actual queda expuesto en rutas concretas:

\begin{itemize}
  \item \code{/claim} y \code{/claim/\{token\}};
  \item \code{/password-reset} y \code{/password-reset/\{token\}};
  \item \code{/login}.
\end{itemize}

En servicios, \code{UserServiceImpl} concentra los m\'etodos \code{issueClaimForEmail}, \code{issuePasswordResetForEmail} y la confirmaci\'on del token. En persistencia se agrega \code{UserPasswordResetToken} y soporte JDBC para emitir, invalidar y consumir tokens.

\section{Hardening de remember-me y cambio de idioma}

El commit \commit{cc284a1d} endurece dos puntos sensibles:

\begin{itemize}
  \item obliga a tener \code{security.rememberme.key} en vez de admitir una key predecible;
  \item hace session-scoped el cambio de idioma, evitando efectos laterales entre requests.
\end{itemize}

El resultado es importante porque corrige un problema de seguridad real y otro de consistencia funcional. No es un simple cleanup.

\section{Registro unificado}

Una de las decisiones de producto del Sprint 2 fue abandonar la separaci\'on de \code{/register/reviewer} y \code{/register/owner}. El commit \commit{4e884b77} materializa ese cambio sobre \code{AuthController} y \code{RegisterForm}.

La implementaci\'on actual funciona as\'i:

\begin{itemize}
  \item \code{/register} sirve un formulario unificado;
  \item el formulario permite elegir el tipo de cuenta;
  \item el controller deriva a \code{userService.registerOwner(...)} o \code{userService.registerReviewer(...)} seg\'un la opci\'on seleccionada;
  \item las vistas viejas quedan desplazadas por una sola pantalla principal.
\end{itemize}

Este punto es central para \story{PAW20-22} porque cambia tanto la superficie web como la forma en que el backend interpreta el alta de usuario.

\section{Verificaci\'on por correo y autologin}

El commit \commit{10915c6a} completa el cierre del alta de cuenta:

\begin{itemize}
  \item agrega estado de verificaci\'on al modelo de usuario;
  \item extiende schema en PostgreSQL y HSQLDB;
  \item agrega el mail \code{register-confirmation.html};
  \item implementa \code{/register/pending}, \code{/register/pending/resend} y \code{/register/confirm/\{token\}};
  \item resuelve autologin luego de confirmar el token.
\end{itemize}

La clave aqu\'i es que la verificaci\'on no qued\'o solo como check booleano en la base. El flujo web se rearm\'o para:

\begin{enumerate}
  \item registrar;
  \item llevar a pantalla intermedia de pendiente;
  \item enviar CTA por correo;
  \item confirmar con token;
  \item volver al sitio con sesi\'on ya abierta.
\end{enumerate}

\subsection{C\'omo se genera el token de confirmaci\'on}

La parte m\'as importante de esta implementaci\'on est\'a en \code{UserServiceImpl}. El flujo real de \code{register(...)} no se limita a crear el usuario y asignar un rol:

\begin{enumerate}
  \item persiste el usuario con \code{passwordHash} ya codificado;
  \item asigna \code{ROLE\_REVIEWER} o \code{ROLE\_RESTAURANT\_OWNER};
  \item llama a \code{createClaimToken(userId)};
  \item env\'ia \code{register-confirmation.html} con la URL \code{/register/confirm/\{token\}};
  \item redirige a \code{/register/pending}.
\end{enumerate}

\code{createClaimToken(...)} implementa una estrategia segura y bastante limpia:

\begin{itemize}
  \item genera un token raw de 32 bytes con \code{SecureRandom};
  \item lo codifica en Base64 URL-safe sin padding;
  \item calcula \code{SHA-256} del token raw;
  \item invalida tokens de claim pendientes previos del mismo usuario;
  \item guarda solamente el hash en base, nunca el token raw;
  \item persiste \code{expires\_at = now + 24 horas}.
\end{itemize}

Eso significa que el enlace enviado por mail contiene el token utilizable, pero la base s\'olo conserva su huella. Si alguien lee la tabla \code{user\_claim\_tokens}, no puede reconstruir el link original. Adem\'as, al invalidar tokens anteriores se evita que un usuario acumule links viejos todav\'ia activos.

\subsection{D\'onde queda persistido y cu\'ando vence}

La informaci\'on de verificaci\'on se reparte entre dos lugares:

\begin{itemize}
  \item \code{users.email\_verified\_at}: marca que la cuenta ya confirm\'o el mail;
  \item \code{user\_claim\_tokens}: guarda \code{token\_hash}, \code{expires\_at}, \code{consumed\_at} y \code{created\_at}.
\end{itemize}

El criterio de validez es expl\'icito:

\begin{itemize}
  \item \code{consumed\_at} debe ser \code{null};
  \item \code{expires\_at} debe ser posterior a \code{LocalDateTime.now()}.
\end{itemize}

En otras palabras, el token dura 24 horas desde su emisi\'on y es de un solo uso. Cuando se confirma la cuenta, el sistema rellena \code{consumed\_at} y ese mismo link deja de servir aunque siga dentro de la ventana temporal.

\subsection{Qu\'e hace exactamente la ruta de confirmaci\'on}

\code{AuthController} resuelve \code{/register/confirm/\{token\}}. La secuencia de confirmaci\'on es:

\begin{enumerate}
  \item el controller delega en \code{userService.confirmRegistration(rawToken)};
  \item el service vuelve a hashear el token recibido;
  \item busca el hash en \code{user\_claim\_tokens};
  \item valida que no est\'e vencido ni consumido;
  \item valida que el usuario exista, tenga credenciales y todav\'ia no est\'e verificado;
  \item actualiza \code{users.email\_verified\_at};
  \item marca \code{consumed\_at};
  \item devuelve el usuario ya verificado.
\end{enumerate}

El paso final no es menor: \code{AuthController.authenticateUser(...)} crea un \code{UsernamePasswordAuthenticationToken}, lo inserta en un \code{SecurityContext} nuevo y guarda ese contexto en la sesi\'on HTTP. Por eso, cuando el usuario toca el link del mail, no s\'olo queda marcado como verificado: adem\'as entra al sitio ya autenticado y redirigido seg\'un su rol por \code{AuthRedirectResolver}.

\subsection{Reenv\'io y simetr\'ia con claim y password reset}

El reenv\'io desde \code{/register/pending/resend} no reusa el token previo. En vez de eso vuelve a ejecutar el mismo mecanismo: genera un token raw nuevo, invalida los anteriores y manda un correo fresco. Eso evita el problema de tener varios links equivalentes circulando al mismo tiempo.

El reset de password implementado en \commit{327aeeb7} sigue exactamente la misma idea:

\begin{itemize}
  \item token raw nuevo;
  \item hash persistido en \code{user\_password\_reset\_tokens};
  \item expiraci\'on de 24 horas;
  \item invalidaci\'on de tokens pendientes previos;
  \item consumo del token al completar \code{/password-reset/\{token\}}.
\end{itemize}

La ventaja es conceptual y de mantenimiento: claim, confirmaci\'on de registro y reset de password terminan apoyados sobre la misma pol\'itica de seguridad, en vez de tener tres soluciones distintas.

\subsection{Evidencia manual incorporada al informe}

Durante el QA manual de esta versi\'on se prob\'o el alta completa de un reviewer nuevo y de un owner nuevo:

\begin{itemize}
  \item alta en \code{/register};
  \item llegada a \code{/register/pending};
  \item confirmaci\'on por link;
  \item autologin posterior a la confirmaci\'on.
\end{itemize}

\qafigure{assets/manual-qa/register-reviewer-pending.png}{Pantalla intermedia de registro pendiente, donde el sistema ya pas\'o del alta local al flujo de verificaci\'on por correo.}

\qafigure{assets/manual-qa/reviewer-confirmed-autologin-home.png}{Resultado de confirmar el token: la cuenta queda verificada y la sesi\'on se abre autom\'aticamente.}

En \code{logs/forkd.log} existe adem\'as evidencia operativa de este flujo: aparecen entradas \emph{Registration confirmed for userId=18} y \emph{Registration confirmed for userId=19}, correspondientes al reviewer y al owner creados durante la prueba manual local.

\section{Cambio de password y borrado de cuenta desde perfil}

La story de ``cambiar password desde perfil'' termina de quedar dentro del producto con \commit{d6ea88ab}. El cambio agrega:

\begin{itemize}
  \item \code{/profile/\{username\}/change-password};
  \item \code{/profile/\{username\}/delete} en GET y POST;
  \item la vista \code{profile/delete.jsp};
  \item soporte DAO y service para borrar reviewer;
  \item textos e i18n asociados.
\end{itemize}

T\'ecnicamente es importante porque desplaza un caso de uso desde ``auth puro'' hacia el perfil del usuario, como se hab\'ia pedido en las notas del Sprint 2.

\chapter{Perfiles, reviews y listas personales}

\section{Perfil propio y perfil ajeno}

Las stories \story{PAW20-14} y \story{PAW20-15} se apoyan sobre \code{ProfileController} y sobre las vistas \code{profile/profile.jsp}. Dentro del rango, los commits \commit{fa2ca772} y \commit{fd9ab781} mejoran el perfil y lo vuelven navegable y paginable:

\begin{itemize}
  \item soporte para \code{/profile/\{profileKey\}};
  \item mejor visualizaci\'on de reviews;
  \item paginaci\'on de reviews;
  \item ordenamiento por criterio.
\end{itemize}

No todo el concepto de perfil nace estrictamente en este rango, pero s\'i queda consolidado y mejorado de forma visible.

\section{Reviews verificadas}

El commit \commit{2e16c230} resuelve una mejora muy visible de producto: distinguir reviews hechas tras una reserva completada. Esto se implementa de punta a punta:

\begin{itemize}
  \item \code{Review} suma el enlace con reserva;
  \item el schema agrega soporte de \code{reservation\_id};
  \item \code{ReviewServiceImpl} valida que solo una reserva completada y propia pueda originar una review verificada;
  \item las vistas \code{restaurant.jsp}, \code{profile/profile.jsp}, \code{profile/reservations.jsp}, \code{owner/preview.jsp} y \code{owner/reviews.jsp} muestran la marca visual.
\end{itemize}

Este punto es relevante porque no fue una decoraci\'on UI aislada: hubo cambio de modelo, persistencia, reglas de negocio y presentaci\'on.

\section{Avatar por defecto y presentaci\'on de reviewer}

El commit \commit{f1dc3686} introduce \code{userAvatar.tag} y reusa la misma representaci\'on visual por defecto en perfil, preview owner y detalle de restaurante. Esto cierra una inconsistencia concreta que exist\'ia en \code{/restaurants/\{id\}}, donde antes se mostraba un c\'irculo gris sin letra.

\section{Restaurantes guardados}

La story \story{PAW20-19} queda soportada por dos commits:

\begin{itemize}
  \item \commit{a89be6e1}, que introduce la lista de guardados;
  \item \commit{5974e49e}, que agrega paginaci\'on y orden.
\end{itemize}

La implementaci\'on se reparte entre:

\begin{itemize}
  \item \code{SavedRestaurantJdbcDao} y \code{SavedRestaurantServiceImpl};
  \item \code{ProfileController};
  \item \code{RestaurantController} para guardar desde el detalle;
  \item la vista \code{profile/saved-restaurants.jsp};
  \item el enum \code{SavedRestaurantSortOrder}.
\end{itemize}

Adem\'as de crear la lista, la versi\'on final agrega:

\begin{itemize}
  \item vac\'io sem\'antico con CTA a explorar;
  \item remoci\'on desde la propia lista;
  \item i18n completo;
  \item tests MVC espec\'ificos.
\end{itemize}

\subsection{C\'omo se implementa el guardado}

El flujo de guardado es simple en UI, pero est\'a bien encapsulado en backend:

\begin{enumerate}
  \item el detalle del restaurante dispara \code{POST /restaurants/\{id\}/save};
  \item \code{RestaurantController} no toca la base directamente: delega en \code{SavedRestaurantServiceImpl};
  \item el service valida dos cosas antes de persistir:
    \begin{itemize}
      \item que el usuario sea reviewer;
      \item que el restaurante sea p\'ublico y visible.
    \end{itemize}
  \item seg\'un el estado actual, llama a \code{save(...)} o \code{remove(...)} sobre \code{SavedRestaurantDao};
  \item el controller vuelve a redirigir al detalle del restaurante para que el bot\'on refleje el nuevo estado.
\end{enumerate}

El detalle importante es que el chequeo de permisos no qued\'o en la vista. Aunque el bot\'on s\'olo se expone en contexto reviewer, el service vuelve a verificar que el usuario y el restaurante sean v\'alidos antes de tocar persistencia.

\subsection{Evidencia manual}

Durante el QA se guard\'o el restaurante \emph{Kansas Palermo} con la cuenta reviewer creada para esta corrida. El estado visual del CTA cambi\'o de \emph{Guardar} a \emph{Quitar de guardados} y la base local qued\'o con la relaci\'on \code{saved\_restaurants(user\_id=18, restaurant\_id=45)}.

\qafigure{assets/manual-qa/reviewer-saved-restaurant.png}{Detalle del restaurante luego de guardarlo con una cuenta reviewer nueva.}

\section{Paginaci\'on configurable en perfiles}

El commit \commit{dc309493} introduce \code{PageSizeResolver} con dos juegos de opciones:

\begin{itemize}
  \item est\'andar: 5, 10, 15, 20;
  \item explore: 9, 15, 21, 27, 30.
\end{itemize}

Este cambio aterriza un pedido expl\'icito de la c\'atedra: evitar listas fijas de 5 elementos cuando el usuario deber\'ia poder elegir el tama\~no de p\'agina. En el perfil repercute sobre reviews, reservas y guardados.

\chapter{Dashboard owner, moderaci\'on admin y ciclo de restaurante}

\section{Dashboard owner y moderaci\'on inicial}

El commit \commit{03bf6161} es el primer gran salto del lado owner/admin dentro del rango:

\begin{itemize}
  \item agrega \code{OwnerDashboardController};
  \item reorganiza \code{AdminController};
  \item suma vistas \code{owner/dashboard.jsp}, \code{owner/preview.jsp}, \code{owner/reviews.jsp};
  \item incorpora la moderaci\'on de restaurantes y el flujo de reviews autenticadas.
\end{itemize}

En t\'erminos de producto, este commit abre la posibilidad de que el owner deje de depender del admin para revisar su propio estado y administrar su superficie operativa.

\subsection{Qu\'e hace concretamente el dashboard owner}

\code{OwnerDashboardController} no es s\'olo una lista. En la versi\'on final concentra tres contratos diferentes:

\begin{itemize}
  \item listado de restaurantes del owner con paginaci\'on y acciones;
  \item alta de un restaurante nuevo en \code{/my-restaurants/new};
  \item edici\'on del restaurante ya existente, adem\'as de sus subflujos de men\'u y reservas.
\end{itemize}

La parte importante es que el owner trabaja sobre el mismo modelo de formulario que el admin, pero con reglas de estado distintas. Cuando el owner crea o modifica datos sensibles del restaurante, la entidad no queda publicada por defecto: vuelve al circuito de moderaci\'on.

\qafigure{assets/manual-qa/owner-dashboard-empty.png}{Dashboard del owner nuevo luego de confirmar la cuenta. El flujo ya lo redirige al panel correcto seg\'un el rol.}

\section{Asignaci\'on de restaurantes hu\'erfanos}

La story de asignar restaurantes sin owner se resuelve con \commit{88e76319}. El flujo se implementa con:

\begin{itemize}
  \item m\'etodos nuevos en \code{RestaurantDao} y \code{UserDao};
  \item soporte service en \code{RestaurantServiceImpl} y \code{UserServiceImpl};
  \item \code{AssignOwnerForm};
  \item validador \code{ExistingOwnerUsername};
  \item vistas \code{admin/assign-owner.jsp} y \code{admin/unassigned-restaurants.jsp}.
\end{itemize}

Esto es un buen ejemplo de feature administrativa que no pod\'ia quedar solo en la vista: necesitaba consultas espec\'ificas, validaci\'on de owner real y actualizaci\'on persistente.

\section{Aprobaci\'on, rechazo y reenv\'io a pending}

Las correcciones pedidas en las notas del Sprint 2 sobre approve/reject se reparten entre \commit{51b7c30b} y \commit{178d5ee5}. El objetivo funcional fue:

\begin{itemize}
  \item sacar al restaurante rechazado de la lista de pendientes;
  \item mostrar el estado rechazado al owner;
  \item permitir que el owner edite y vuelva a enviar a pending;
  \item notificar al owner por correo en su idioma.
\end{itemize}

La implementaci\'on real se concentra en \code{RestaurantServiceImpl.approve(...)} y \code{RestaurantServiceImpl.reject(...)}. Ah\'i se actualiza estado, publicaci\'on y mail. En \code{logs/forkd.log} aparecen entradas de negocio como \emph{Restaurant rejected restaurantId=31}, lo que confirma uso real del flujo.

\subsection{C\'omo se model\'o el ciclo de estados}

El ciclo de aprobaci\'on no es un simple flag visual. En \code{RestaurantServiceImpl} se ve una m\'aquina de estados razonablemente clara:

\begin{itemize}
  \item \code{submitRestaurant(...)} crea el restaurante, reconstruye el layout de mesas, asigna owner, deja el estado en \code{PENDING} y fuerza \code{unpublish()};
  \item \code{approve(...)} hace \code{updateStatus(APPROVED)} y luego \code{publish()};
  \item \code{reject(...)} hace \code{updateStatus(REJECTED, reason, reviewedAt)} y luego \code{unpublish()};
  \item \code{updateOwnedRestaurant(...)} vuelve a \code{PENDING} si el restaurante estaba rechazado o si un owner aprob\'o cambios sensibles sobre nombre, direcci\'on, barrio, cocina o rango de precio.
\end{itemize}

Eso implementa exactamente la regla funcional buscada: el owner puede operar sobre su restaurante, pero no puede autopublicar modificaciones sensibles. El admin sigue siendo quien cierra el circuito.

\subsection{Wizard de alta/edici\'on}

La separaci\'on posterior en \commit{28014edf} agrega un wizard con cuatro pasos:

\begin{enumerate}
  \item identidad;
  \item ubicaci\'on;
  \item presentaci\'on;
  \item horarios.
\end{enumerate}

Ese wizard no es puramente cosm\'etico. El backend recibe finalmente un \code{RestaurantForm} unificado y lo traduce a \code{SaveRestaurantData}; luego el service decide si est\'a creando, actualizando, manteniendo assets previos o forzando reenv\'io a moderaci\'on.

\qafigure{assets/manual-qa/owner-restaurant-wizard-hours.png}{Paso final del wizard owner para horarios. La edici\'on ya no vive en un formulario monol\'itico.}

\subsection{Evidencia manual de admin}

Durante el QA manual se ingres\'o tambi\'en con \code{admin\_forkd} y se ejecut\'o una aprobaci\'on real sobre un restaurante pendiente del entorno local. El flujo recorrido fue:

\begin{enumerate}
  \item login en \code{/login};
  \item entrada al panel \code{/admin/restaurants};
  \item acceso a \code{/admin/pending};
  \item aprobaci\'on efectiva de la solicitud pendiente;
  \item retorno al listado sin pendientes y con toast de confirmaci\'on.
\end{enumerate}

\qafigure{assets/manual-qa/admin-pending-approval.png}{Panel admin antes de aprobar un restaurante pendiente.}

\qafigure{assets/manual-qa/admin-pending-approved.png}{Resultado luego de aprobar el restaurante pendiente desde el backoffice admin.}

\section{Preview del restaurante desde owner}

La story \story{PAW20-55} se cerr\'o con \commit{26872714}. El problema no era meramente est\'etico: el bot\'on de preview apuntaba a una ruta owner, pero el preview correcto deb\'ia abrir la vista p\'ublica o privada visible del restaurante.

La correcci\'on se hizo en dos pasos:

\begin{itemize}
  \item \code{owner/dashboard.jsp} pasa a enlazar a \code{/restaurants/\{id\}};
  \item \code{RestaurantController} agrega \code{findVisibleRestaurant(...)} para permitir que owner o admin vean restaurantes no publicados, sin abrir esa capacidad a usuarios an\'onimos.
\end{itemize}

Es un cambio chico en superficie, pero correcto en dise\~no porque corrige routing y permisos al mismo tiempo.

\section{Separaci\'on de edici\'on: restaurante, men\'u y reservas}

La nota del Sprint 2 ped\'ia no editar men\'u y reservas dentro del mismo formulario gigante. Eso se traduce en:

\begin{itemize}
  \item \commit{2e52cd5a}, que separa formalmente men\'u, configuraci\'on de reservas y datos generales;
  \item \commit{28014edf}, que agrega un wizard para la edici\'on/creaci\'on de restaurante.
\end{itemize}

T\'ecnicamente este cambio fue grande porque oblig\'o a partir el contrato del formulario en varias piezas:

\begin{itemize}
  \item \code{RestaurantForm};
  \item \code{MenuForm};
  \item \code{ReservationConfigForm};
  \item fragmentos compartidos \code{restaurant-form-body.jspf}, \code{menu-form-body.jspf} y \code{reservation-config-form-body.jspf};
  \item nuevas rutas GET/POST para \code{/my-restaurants/\{id\}/menu} y \code{/my-restaurants/\{id\}/reservation-settings}, adem\'as de sus equivalentes admin.
\end{itemize}

La consecuencia es que el owner ya no actualiza ``todo el restaurante'' para tocar solo el men\'u o solo la pol\'itica de reservas.

\section{Men\'u enriquecido con descripci\'on e imagen por plato}

La story \story{PAW20-28} recibe una mejora fuerte con \commit{944a5e9a}. El men\'u deja de ser solo texto y pasa a modelar:

\begin{itemize}
  \item secciones;
  \item items con nombre, descripci\'on e imagen;
  \item renderizaci\'on m\'as rica en la vista \code{restaurantMenu.jsp};
  \item edici\'on m\'as cuidada desde \code{owner/menu-form.jsp}.
\end{itemize}

Esto tambi\'en explica por qu\'e el men\'u tuvo que separarse del resto del formulario: con m\'as estructura, ya no era sostenible como simple bloque embebido.

\section{Otras mejoras sobre el ciclo de restaurante}

En este rango tambi\'en entran varios ajustes relacionados:

\begin{itemize}
  \item \commit{bff8d394}: agrega la direcci\'on al encabezado del detalle del restaurante;
  \item \commit{0e675289}: unifica rutas en ingl\'es;
  \item \commit{78445ded}: resuelve drift entre tags legacy y tags reales mediante normalizaci\'on de cat\'alogo y migraci\'on de limpieza;
  \item \commit{72a5d7d8}: ajusta longitud m\'axima de direcci\'on del restaurante;
  \item \commit{6c41ef92}, \commit{bd459157}, \commit{036911dd} y \commit{342ac6cc}: corrigen detalles de padding y presentaci\'on en formularios y listados.
\end{itemize}

\chapter{Sistema de reservas}

\section{Visi\'on general}

El sistema de reservas es la implementaci\'on m\'as grande del rango. Se materializa principalmente a trav\'es del merge \commit{5589cdb9}, pero conviene explicarlo por fases porque as\'i fue dise\~nado y as\'i lo documenta la wiki.

\section{Fase 1: modelo, schema y persistencia base}

El commit \commit{37e7997e} agrega la base del subsistema:

\begin{itemize}
  \item enums de dominio como \code{AcceptsReservationsStatus}, \code{ReservationConfirmationMode}, \code{ReservationDisabledMode}, \code{ReservationSlotDuration} y \code{ReservationStatus};
  \item entidades \code{Reservation}, \code{RestaurantTable} y \code{TimeSlot};
  \item contratos \code{ReservationDao} y \code{RestaurantTableDao};
  \item implementaciones \code{ReservationJdbcDao} y \code{RestaurantTableJdbcDao};
  \item extensiones de \code{Restaurant} y \code{Review};
  \item cambios de schema en \code{schema-base.sql}, \code{schema-postgres.sql} y \code{schema-hsqldb.sql};
  \item fixtures y tests JDBC.
\end{itemize}

No es una fase menor: define el vocabulario del problema y crea la estructura f\'isica sobre la que luego se apoyan slots, estados, notificaciones y reviews verificadas.

\section{Fase 2: configuraci\'on de reservas dentro del restaurante}

El commit \commit{5151f2f5} introduce el subform de configuraci\'on de reservas y la regla de ``layout congelado'':

\begin{itemize}
  \item \code{ReservationConfigForm};
  \item \code{RestaurantTableGroupForm};
  \item \code{ReservationConfigValidator};
  \item \code{ValidReservationConfig};
  \item \code{LayoutFrozenException};
  \item extensi\'on de \code{SaveRestaurantData}.
\end{itemize}

Esto resuelve el caso de uso \story{PAW20-31}: cu\'antas mesas disponibles tiene el restaurante y con qu\'e layout opera.

\section{Fase 3: disponibilidad y asignaci\'on autom\'atica}

El commit \commit{2567832d} agrega la capa algor\'itmica del sistema:

\begin{itemize}
  \item \code{SlotCatalog} para derivar slots a partir de horarios y duraci\'on;
  \item l\'ogica de asignaci\'on de mesa seg\'un cantidad de personas;
  \item consultas JDBC para detectar disponibilidad por slot;
  \item tests unitarios espec\'ificos del algoritmo.
\end{itemize}

La decisi\'on relevante de arquitectura fue hacer la resoluci\'on principal en services y dejar a persistence el soporte de datos necesarios, en vez de dispersar la l\'ogica en JSP o controllers.

\subsection{C\'omo se calculan los slots}

\code{SlotCatalog} toma tres entradas:

\begin{itemize}
  \item la lista de \code{RestaurantHour};
  \item la duraci\'on del slot;
  \item la fecha a consultar.
\end{itemize}

Primero filtra los horarios del d\'ia de semana pedido. Luego, para cada franja, genera slots desde \code{openTime} hasta \code{closeTime} sumando la duraci\'on configurada. Si el \'ultimo slot no entra completo, recorta el final al horario de cierre. Si detecta turnos overnight o datos inconsistentes, los descarta y deja warning en logs.

Despu\'es, \code{ReservationServiceImpl.listAvailableSlots(...)} cruza ese cat\'alogo te\'orico contra las mesas realmente libres en ese horario. El resultado visible al usuario no es ``todos los slots del horario'', sino s\'olo los que siguen teniendo al menos una mesa compatible con la cantidad de personas.

\section{Fase 4: creaci\'on de reservas desde el detalle del restaurante}

Con \commit{03fd92b8} aparece la superficie visible para el reviewer:

\begin{itemize}
  \item \code{ReservationController};
  \item \code{/restaurants/\{id\}/reservations};
  \item \code{/restaurants/\{id\}/reservations/slots} para poblar slots;
  \item \code{CreateReservationForm};
  \item \code{CreateReservationFormValidator};
  \item \code{ReservationPageSupport} y \code{RestaurantDetailViewSupport};
  \item la tabla de notificaciones \code{reservation\_notifications};
  \item excepciones \code{InvalidReservationSlotException}, \code{SlotAlreadyPastException}, \code{SlotUnavailableException} y \code{ReservationsDisabledException}.
\end{itemize}

El commit de ajuste \commit{fed1b8d8} corrige detalles de estas primeras fases, especialmente en advice, fixtures y tests.

\subsection{C\'omo fluye una reserva de punta a punta}

La implementaci\'on final de la reserva tiene varias validaciones encadenadas:

\begin{enumerate}
  \item el reviewer selecciona fecha y cantidad de personas;
  \item el frontend de \code{restaurant.jsp} consulta \code{/restaurants/\{id\}/reservations/slots};
  \item el backend valida que el restaurante tenga reservas habilitadas, duraci\'on de slot y horarios definidos;
  \item \code{SlotCatalog} genera el cat\'alogo posible para ese d\'ia;
  \item \code{RestaurantTableDao} devuelve qu\'e comienzos de slot siguen teniendo capacidad libre;
  \item el usuario elige uno de esos slots y confirma;
  \item \code{createReservation(...)} revalida todo, busca la primera mesa compatible y persiste la reserva.
\end{enumerate}

La elecci\'on del estado inicial tambi\'en est\'a codificada en service y no en la vista:

\begin{itemize}
  \item si el restaurante usa confirmaci\'on manual, la reserva nace \code{PENDING};
  \item si usa confirmaci\'on autom\'atica, el sistema confirma de inmediato salvo cuando la mesa asignada es demasiado grande en relaci\'on con el grupo, caso en el que fuerza revisi\'on manual.
\end{itemize}

Durante el QA manual local se prob\'o este flujo con la cuenta \emph{Reviewer QA}, creando una reserva real en \emph{Kansas Palermo}. La reserva qued\'o persistida como \code{reservationId=12}.

\qafigure{assets/manual-qa/reviewer-reservation-created.png}{Reserva creada desde la p\'agina del restaurante y visible luego en \emph{Mis Reservas}.}

\section{Fase 5: gesti\'on de reservas por parte del usuario}

La story \story{PAW20-24} no termina con crear reservas; tambi\'en hace falta verlas y cancelarlas. Eso llega con \commit{b99dcafb} y se pule en \commit{175f90bf}.

La implementaci\'on actual incluye:

\begin{itemize}
  \item \code{ProfileController} sobre \code{/profile/\{profileKey\}/reservations};
  \item \code{ReservationServiceImpl.listUserReservations(...)} y \code{cancelAsUser(...)};
  \item ordenamiento por \code{ReservationSort};
  \item cancelaci\'on con actualizaci\'on de estado y notificaci\'on;
  \item vista \code{profile/reservations.jsp}.
\end{itemize}

Esto tambi\'en refuerza el valor del perfil como centro operativo del reviewer.

\section{Fase 6: gesti\'on owner de reservas}

La fase owner aparece en \commit{5387190a}. Es una de las partes m\'as interesantes del dise\~no porque mezcla permisos, negocio y UX:

\begin{itemize}
  \item \code{OwnerReservationController} atiende \code{/my-restaurants/\{id\}/reservations};
  \item \code{ReservationServiceImpl} expone \code{listReservationsByRestaurant(...)}, \code{listReassignableTables(...)}, \code{findPageContainingReservation(...)}, \code{confirmAsOwner(...)} y \code{cancelAsOwner(...)};
  \item \code{ReservationStatusFilter} y \code{ReservationSort} soportan filtros de estado y orden;
  \item el owner puede confirmar pendientes, cancelarlas y reasignar una mesa alternativa si sigue cumpliendo la capacidad requerida;
  \item la vista \code{owner/reservations.jsp} conserva el contexto mediante \code{page}, \code{sort} y \code{focus}.
\end{itemize}

Desde el punto de vista t\'ecnico, esta fase est\'a bien resuelta porque no arrastra la lista completa a memoria para resolver la paginaci\'on de foco. El servicio computa la p\'agina que contiene una reserva dada a trav\'es de SQL, preservando el orden actual.

\subsection{Confirmaci\'on, cancelaci\'on y reasignaci\'on}

Cuando el owner confirma una reserva pendiente, \code{ReservationServiceImpl.confirmAsOwner(...)} hace cuatro cosas importantes:

\begin{enumerate}
  \item carga la reserva y exige que su estado sea \code{PENDING};
  \item opcionalmente permite reasignar una mesa, pero s\'olo entre las que siguen siendo v\'alidas para ese mismo slot;
  \item actualiza el estado a \code{CONFIRMED};
  \item dispara el mail de confirmaci\'on si todav\'ia no se hab\'ia emitido uno de ese tipo.
\end{enumerate}

La cancelaci\'on owner sigue la misma l\'ogica de negocio: s\'olo se puede cancelar si la reserva estaba \code{PENDING} o \code{CONFIRMED}, se actualiza a \code{CANCELLED} y se dispara el correo correspondiente al reviewer.

Durante el QA se ingres\'o con el owner seed \emph{kansas}, se abri\'o \code{/my-restaurants/45/reservations} y se confirm\'o manualmente la reserva reci\'en creada por el reviewer nuevo.

\qafigure{assets/manual-qa/owner-reservations-pending.png}{Panel owner de reservas con una reserva pendiente, mesa reasignable y acciones de confirmaci\'on o rechazo.}

\section{Fase 7: scheduler de expiraci\'on y completado}

La automatizaci\'on aparece con \commit{155ab1b2}:

\begin{itemize}
  \item se agrega \code{BackgroundExecutionConfig};
  \item \code{ReservationSchedulerImpl} expone \code{tick()} y \code{sendDailyReminders()};
  \item el scheduler borra pendientes expiradas y marca como completadas reservas confirmadas cuyo slot ya termin\'o;
  \item el setup se integra con \code{ServicesConfig} y \code{web.xml};
  \item se agregan tests dedicados.
\end{itemize}

La evidencia operativa aparece en los logs actuales: \code{logs/forkd.log} contiene decenas de l\'ineas con \emph{Reservation scheduler tick completed: deletedPending=0 completedConfirmed=0}, lo que prueba ejecuci\'on real del proceso background.

\section{Fase 8: mails transaccionales y recordatorio diario}

Con \commit{3b6c5221} el sistema pasa de gestionar estados a comunicar esos estados. Se agregan:

\begin{itemize}
  \item templates \code{reservation-confirmed.html}, \code{reservation-cancelled.html}, \code{reservation-reminder.html}, \code{reservation-completed.html} y \code{reservation-pending-owner.html};
  \item layout com\'un \code{_layout.html};
  \item m\'etodos de \code{SmtpMailServiceImpl} para cada notificaci\'on;
  \item disparo de mails desde services y scheduler, no desde webapp;
  \item idempotencia apoyada sobre \code{reservation\_notifications}.
\end{itemize}

Este punto tambi\'en cubre parcialmente \story{PAW20-32} y \story{PAW20-34}: el usuario recibe mail cuando se cancela una reserva, y varios correos pasan a contener CTAs que reingresan al flujo web adecuado.

\section{Fase 9: reviews verificadas ligadas a reservas}

Aunque ya se explicaron dentro del perfil, funcionalmente esta fase pertenece a reservas. \commit{2e16c230} cierra el puente entre:

\begin{itemize}
  \item una reserva completada;
  \item la habilitaci\'on para dejar review;
  \item la marca visual de review verificada.
\end{itemize}

El dise\~no evita que la verificaci\'on sea una etiqueta arbitraria de UI: se deriva de una relaci\'on persistida y validada en backend.

\subsection{C\'omo se habilita la rese\~na verificada}

La ruta normal del detalle del restaurante no habilita sola la review verificada. El mecanismo completo es este:

\begin{enumerate}
  \item debe existir una reserva \code{COMPLETED};
  \item esa reserva debe pertenecer al usuario autenticado;
  \item la reserva debe corresponder al mismo restaurante;
  \item el detalle se abre con \code{/restaurants/\{id\}?review=\{reservationId\}};
  \item \code{RestaurantController.detail(...)} consulta \code{reviewService.canReviewReservation(...)} y, si corresponde, inyecta \code{reviewReservationId} en el modelo;
  \item el formulario renderiza un hidden \code{reservationId} y una marca visual de ``rese\~na verificada'' antes de publicar.
\end{enumerate}

Al enviar la review, \code{ReviewServiceImpl.create(...)} vuelve a validar que la reserva:

\begin{itemize}
  \item exista;
  \item pertenezca al usuario;
  \item pertenezca al restaurante;
  \item est\'e en estado \code{COMPLETED};
  \item no tenga ya otra review asociada.
\end{itemize}

Si cualquiera de esas condiciones falla, la review no se crea como verificada. Reci\'en cuando todas pasan, se persiste \code{reservation\_id} en la tabla de reviews y la UI puede mostrar el badge \emph{Verificada} con respaldo real de datos.

\subsection{Evidencia manual}

Para poder cerrar el ciclo durante la prueba manual local, la reserva creada en QA se marc\'o como completada en la base local y luego se abri\'o \code{/restaurants/45?review=12}. El formulario mostr\'o el badge de rese\~na verificada antes de publicar y la review final qued\'o persistida con \code{reservation\_id=12}.

\qafigure{assets/manual-qa/reviewer-verified-review-form.png}{Formulario de review ya enlazado a una reserva completada, mostrando la marca de rese\~na verificada antes de publicar.}

\qafigure{assets/manual-qa/reviewer-verified-review-published.png}{Review publicada y visible en el detalle del restaurante con badge de verificaci\'on.}

\section{Fase 10: hardening final}

El commit \commit{778065ed} termina de consolidar el subsistema:

\begin{itemize}
  \item i18n sim\'etrico;
  \item mapeos de error y advice;
  \item cierres de seguridad;
  \item actualizaci\'on del README;
  \item ajustes en \code{ExploreController}, \code{OwnerDashboardController} y \code{ProfileController}.
\end{itemize}

La wiki considera esta fase el cierre operativo del sistema de reservas, y el documento \code{2026-04-17_cierre-implementacion-reservas_v1.md} lo da por \emph{production-ready} dentro del est\'andar del proyecto.

\section{Problemas encontrados durante la implementaci\'on}

El log del 17 de abril muestra un detalle valioso de la evoluci\'on real: durante el alta del soporte de reviews verificadas hubo un problema de SQL al intentar ejecutar \emph{ALTER TABLE reviews ADD CONSTRAINT IF NOT EXISTS ...} en PostgreSQL. El estado final del repo refleja la correcci\'on:

\begin{itemize}
  \item en \code{schema-postgres.sql} se usa \code{CREATE UNIQUE INDEX IF NOT EXISTS reviews\_reservation\_id\_unique ON reviews(reservation\_id)};
  \item en \code{schema-base.sql} y \code{schema-hsqldb.sql} se mantiene la restricci\'on compatible con el entorno correspondiente.
\end{itemize}

Esto es importante para el informe porque muestra que no fue una implementaci\'on lineal: hubo una fase de estabilizaci\'on real sobre el runtime.

\chapter{Mailing, logs y evidencia de ejecuci\'on}

\section{Mails de negocio fuera del sistema de reservas}

El rango no s\'olo mejor\'o mails de reservas. Tambi\'en refuerza otros dos flujos:

\begin{itemize}
  \item \commit{178d5ee5}: rechazo de restaurante con template \code{restaurant-rejection.html};
  \item \commit{ac0220ff}: CTA en todos los correos, incluyendo \code{new-review.html}.
\end{itemize}

Esto conecta con \story{PAW20-33}: el restaurante recibe notificaci\'on cuando entra una rese\~na, y ya no como email plano o incompleto, sino con CTA navegable hacia la p\'agina correspondiente.

\section{Uso del locale del destinatario}

El commit \commit{07f980da} corrige una deuda importante: los mails a owners y reviewers deben salir en el idioma del destinatario, no en el idioma del request que dispar\'o la acci\'on.

La implementaci\'on refuerza una regla de arquitectura del proyecto:

\begin{itemize}
  \item persistence debe proveer email y locale del due\~no;
  \item services debe decidir y disparar el mail;
  \item webapp no debe contaminar la composici\'on del correo.
\end{itemize}

\section{Evidencia en logs}

Las siguientes entradas observadas en los logs locales son especialmente \'utiles para el informe:

\begin{itemize}
  \item confirmaci\'on de registro: \emph{Registration confirmed for userId=17} en \code{logs/forkd.log};
  \item confirmaci\'on owner de reserva: \emph{Reservation 10 confirmed by owner} en \code{logs/forkd.2026-04-17.log};
  \item rechazo de restaurante: \emph{Restaurant rejected restaurantId=31} en \code{logs/forkd.log};
  \item ejecuci\'on repetida del scheduler: \emph{Reservation scheduler tick completed: deletedPending=0 completedConfirmed=0} en \code{logs/forkd.log}.
\end{itemize}

Estas entradas ayudan a justificar que el informe no se basa s\'olo en tests o en diff est\'atico, sino tambi\'en en ejecuci\'on real del sistema.

\chapter{Explore, paginaci\'on y experiencia de listas}

\section{Persistencia mutua entre b\'usqueda y filtros}

La story \story{PAW20-53} queda claramente respaldada por \commit{e7ba605a}. El problema era que al escribir una b\'usqueda se perd\'ian filtros, o viceversa.

La soluci\'on fue concreta y server-side:

\begin{itemize}
  \item \code{header.tag} empez\'o a reenviar como hidden inputs el estado actual de barrio, cocinas, tags, priceRanges, openNow y sort;
  \item \code{explore.jsp} empez\'o a reenviar la query actual desde el propio formulario lateral;
  \item el usuario ya no necesita recomponer el estado de explore cada vez que toca el search box.
\end{itemize}

\section{Tama\~no de p\'agina fijo pero configurable}

La story sobre \code{pageSize} se materializa en \commit{dc309493}. El cambio agrega un resolvedor expl\'icito:

\begin{itemize}
  \item \code{PageSizeResolver.resolveStandardPageSize(...)} para listas cl\'asicas;
  \item \code{PageSizeResolver.resolveExplorePageSize(...)} para la grilla de explore.
\end{itemize}

La decisi\'on fue deliberadamente conservadora: el usuario puede elegir el tama\~no, pero entre opciones cerradas definidas por el sistema. Eso evita valores arbitrarios y simplifica validaci\'on, performance e integraci\'on UI.

\section{Controles de paginaci\'on m\'as ricos}

La story \story{PAW20-50} se ve bien en \commit{63090c20}. El tag \code{paging.tag} pasa a tener:

\begin{itemize}
  \item bot\'on de primero;
  \item bot\'on de anterior;
  \item ventana de p\'aginas alrededor de la actual;
  \item bot\'on de siguiente;
  \item bot\'on de \'ultima;
  \item labels accesibles e i18n para cada control.
\end{itemize}

Esto corrige la limitaci\'on visual previa y tambi\'en mejora accesibilidad, porque los controles reciben \code{aria-label} y \code{title} expl\'icitos.

\section{Sobre las stories de filtros, orden y rango de precios}

Las capturas del tablero muestran \story{PAW20-36}, \story{PAW20-44} y \story{PAW20-51}. Dentro del rango analizado se puede afirmar con confianza que hubo:

\begin{itemize}
  \item persistencia del estado entre query y filtros;
  \item page size configurable;
  \item ajuste visual de paginaci\'on;
  \item orden y filtros tambi\'en presentes en reservas owner y listas de perfil.
\end{itemize}

Sin embargo, la introducci\'on original de ciertos filtros de explore, incluido el uso de price range como concepto, ya tiene antecedentes antes de \commit{f81eb49}. Por eso, para el informe final conviene presentar esas stories como \emph{refinadas y consolidadas} en este rango, no necesariamente como creadas desde cero aqu\'i.

\chapter{Hardening transversal y limpieza arquitect\'onica}

\section{Control de acceso declarativo}

El commit \commit{6fa58cd5} migra restricciones hacia un modelo m\'as declarativo usando \code{@PreAuthorize} y helpers de acceso. El valor de este cambio es doble:

\begin{itemize}
  \item reduce l\'ogica manual de ownership en controllers;
  \item centraliza decisiones de seguridad en puntos auditables.
\end{itemize}

\section{Validaci\'on declarativa e inline}

La pareja \commit{bcec15de} + \commit{94663e93} vuelve a empujar la validaci\'on hacia formularios y validadores:

\begin{itemize}
  \item \code{ValidUploadedImage};
  \item \code{UniqueProfileEditUsername};
  \item reglas compartidas de \code{UploadedImageRules};
  \item menos l\'ogica de negocio en controllers;
  \item conservaci\'on de errores inline.
\end{itemize}

Este punto es importante porque responde a una preferencia fuerte del proyecto: los errores recuperables deben volver al form, no tirar al usuario a una pantalla global de error.

\section{L\'imites de input y request length}

Los commits \commit{46dd99e} y \commit{931ee306} endurecen longitudes m\'aximas en toda la superficie web:

\begin{itemize}
  \item \code{InputMaxLengths};
  \item \code{RequestLengthInterceptor};
  \item \code{LoginRequestLengthFilter};
  \item \code{RequestParamTooLongException};
  \item nueva vista \code{error/414.jsp}.
\end{itemize}

Ac\'a el hardening no se limit\'o a agregar \code{@Size}: tambi\'en se protegi\'o el request antes de que termine en un estado ambiguo en la capa web.

\section{Correcciones de persistencia y mapeo}

Otros cambios transversales relevantes del rango:

\begin{itemize}
  \item \commit{201bae1f}: elimina N+1 en carga de roles dentro de \code{UserJdbcDao};
  \item \commit{76048cae}: agrega \code{ON DELETE CASCADE} a \code{reviews.user\_id};
  \item \commit{b0a6e5ed}: mueve ownership de configuraci\'on fuera de \code{webapp};
  \item \commit{3e00fcb9}: remueve Java est\'atico embebido;
  \item \commit{120e0432}: vuelve a sacar l\'ogica de negocio fuera de controllers tras el crecimiento del sprint.
\end{itemize}

\section{Errores, logging y packaging}

Finalmente, el rango consolida varios cierres de calidad:

\begin{itemize}
  \item \commit{8cc69875}: controlador y vistas de error coherentes con Spring Security;
  \item \commit{d6e93195}: logging expl\'icito en services y migraciones;
  \item \commit{2da7cd24}: cleanup de packaging;
  \item \commit{59fe8097}, \commit{b54ad5de} y \commit{b01a321a}: estabilizaci\'on incremental de tests.
\end{itemize}

\chapter{QA manual ejecutado para esta versi\'on del informe}

\section{Objetivo}

Adem\'as de releer c\'odigo, wiki y logs, esta versi\'on del informe ejecut\'o pruebas manuales sobre la app local levantada en \code{http://localhost:8080}. El objetivo fue verificar que las historias m\'as importantes del rango no s\'olo existen en el c\'odigo, sino que efectivamente recorren sus rutas completas de UI, service y persistencia.

\section{Usuarios y contexto usados}

Para no depender de usuarios seed ya existentes, se crearon dos cuentas nuevas:

\begin{itemize}
  \item reviewer: \code{reviewer\_0419223402};
  \item owner: \code{owner\_0419223402}.
\end{itemize}

Adicionalmente se usaron dos cuentas seed para validar superficies ya pobladas del sistema:

\begin{itemize}
  \item owner seed: \code{kansas};
  \item admin seed: \code{admin\_forkd}.
\end{itemize}

\section{Resultado de los checks}

\begin{longtable}{p{0.32\textwidth}p{0.47\textwidth}p{0.13\textwidth}}
\toprule
\textbf{Flujo probado} & \textbf{Evidencia obtenida} & \textbf{Resultado} \\
\midrule
\endhead
Registro reviewer nuevo & alta en \code{/register}, pantalla \code{/register/pending}, confirmaci\'on por token y autologin & OK \\
Registro owner nuevo & alta en \code{/register}, confirmaci\'on por token y entrada al dashboard owner vac\'io & OK \\
Guardar restaurante & cambio de CTA a \emph{Quitar de guardados} y fila en \code{saved\_restaurants} & OK \\
Crear reserva reviewer & reserva real en \code{Kansas Palermo}, visible en \code{/profile/\{user\}/reservations} como \code{PENDING} & OK \\
Confirmaci\'on owner & confirmaci\'on desde \code{/my-restaurants/45/reservations} con la cuenta owner del restaurante & OK \\
Review verificada & review publicada con badge \emph{Verificada} y \code{reservation\_id=12} persistido en base & OK \\
Moderaci\'on admin & aprobaci\'on real de un restaurante pendiente desde \code{/admin/pending} & OK \\
\bottomrule
\end{longtable}

\section{Aclaraci\'on sobre el cierre de la reserva}

Para llegar a la rese\~na verificada durante la misma corrida manual fue necesario mover localmente la reserva creada a estado \code{COMPLETED}. Esa intervenci\'on no forma parte del producto final; se hizo s\'olo para destrabar el caso de prueba dentro de la misma sesi\'on de QA, sin esperar la transici\'on natural por fecha y scheduler.

\section{Evidencia t\'ecnica complementaria}

Los resultados del QA tambi\'en quedaron reflejados en datos y logs:

\begin{itemize}
  \item \code{saved\_restaurants} contiene la relaci\'on \code{(user\_id=18, restaurant\_id=45)};
  \item la reserva \code{id=12} qued\'o asociada al reviewer nuevo y luego se us\'o para la review verificada;
  \item la review creada en QA qued\'o persistida con \code{reservation\_id=12};
  \item \code{logs/forkd.log} registra, entre otras l\'ineas, \emph{Reservation created reservationId=12 ... status=PENDING}, \emph{Reservation 12 confirmed by owner} y \emph{Restaurant approved restaurantId=31}.
\end{itemize}

\chapter{Matriz de stories del tablero capturado}

\section{Lectura de la matriz}

La siguiente tabla cruza las tarjetas visibles en las capturas con la evidencia hallada en el rango. Se usa la siguiente escala:

\begin{itemize}
  \item \confidence{Alta}: la feature o fix nace o se cierra claramente dentro del rango;
  \item \confidence{Media}: el rango completa o reorganiza una capacidad existente;
  \item \confidence{Baja}: la tarjeta existe en el tablero, pero el repositorio sugiere que su origen es anterior y aqu\'i s\'olo hay refinamientos.
\end{itemize}

\small
\begin{longtable}{p{0.14\textwidth}p{0.31\textwidth}p{0.22\textwidth}p{0.23\textwidth}}
\toprule
\textbf{Story} & \textbf{Tarea visible} & \textbf{Commits principales} & \textbf{Confianza} \\
\midrule
\endhead
\story{PAW20-22} & El usuario quiere registrarse & \commit{f497af22}, \commit{4e884b77}, \commit{10915c6a} & Alta \\
\story{PAW20-19} & Guardar restaurantes para visitar & \commit{a89be6e1}, \commit{5974e49e} & Alta \\
\story{PAW20-15} & Ver perfiles de otros & \commit{fa2ca772}, \commit{fd9ab781} & Media \\
\story{PAW20-14} & Ver su perfil & \commit{fa2ca772}, \commit{d6ea88ab} & Media \\
\story{PAW20-24} & Hacer una reserva & \commit{37e7997e} a \commit{03fd92b8} & Alta \\
\story{PAW20-27} & Ver todos sus restaurants y editarlos & \commit{03bf6161}, \commit{28014edf}, \commit{2e52cd5a} & Alta \\
\story{PAW20-28} & Agregar su menu al perfil & \commit{944a5e9a}, \commit{2e52cd5a} & Media \\
\story{PAW20-29} & Admin valida restaurants & \commit{03bf6161}, \commit{51b7c30b}, \commit{178d5ee5} & Alta \\
\story{PAW20-30} & El restaurant puede cancelar una reserva & \commit{5387190a}, \commit{3b6c5221} & Alta \\
\story{PAW20-31} & El restaurant define cuantas mesas tiene & \commit{5151f2f5} & Alta \\
\story{PAW20-32} & El usuario recibe un mail si se cancela la reserva & \commit{3b6c5221} & Alta \\
\story{PAW20-33} & El restaurant recibe notificaci\'on por rese\~na & \commit{03bf6161}, \commit{ac0220ff} & Media \\
\story{PAW20-34} & Acceder a la pagina desde el mail & \commit{3b6c5221}, \commit{10915c6a}, \commit{ac0220ff} & Media \\
\story{PAW20-35} & Paginaci\'on de ancho fijo & \commit{63090c20} & Alta \\
\story{PAW20-36} & Ordenar junto a filtros & antecedentes previos + ajustes en \commit{e7ba605a} & Baja \\
\story{PAW20-44} & Buscar, filtrar y ordenar al mismo tiempo & antecedentes previos + ajustes en \commit{e7ba605a} & Baja \\
\story{PAW20-50} & Bot\'on al primero y \'ultimo & \commit{63090c20} & Alta \\
\story{PAW20-51} & Indicar rango de precios & antecedentes previos del explore; consolidaci\'on en rango & Baja \\
\story{PAW20-53} & Al buscar se pierde la consulta & \commit{e7ba605a} & Alta \\
\story{PAW20-55} & Arreglar bot\'on preview en /mis-restaurantes & \commit{26872714} & Alta \\
\bottomrule
\end{longtable}
\normalsize

\chapter{Conclusiones}

\section{Balance}

El rango posterior a \commit{f81eb49} no fue una simple fase de ``pulido''. De hecho, all\'i se termina de construir gran parte del producto operativo:

\begin{itemize}
  \item se cierra el flujo completo de cuenta de usuario;
  \item se crea y endurece el sistema de reservas;
  \item se profesionaliza el mailing;
  \item se vuelve utilizable el dashboard owner y la moderaci\'on admin;
  \item se refuerzan validaci\'on, seguridad, i18n y arquitectura.
\end{itemize}

\section{Checks manuales que siguen siendo relevantes}

La propia nota de Sprint 2 deja dos seguimientos operativos abiertos:

\begin{itemize}
  \item prueba manual integral del flujo de reservas;
  \item prueba manual integral del flujo auth.
\end{itemize}

Eso no invalida la implementaci\'on; simplemente marca que el cierre del sprint mezcla c\'odigo ya integrado con una capa de QA manual todav\'ia necesaria.

\section{Siguiente versi\'on sugerida del informe}

Para una v2 del documento conviene agregar:

\begin{itemize}
  \item capturas exportadas del tablero y de la aplicaci\'on ya embebidas en el PDF;
  \item una cronolog\'ia compacta por fecha;
  \item un ap\'endice con las rutas finales del sistema y los formularios principales;
  \item una tabla adicional \emph{story $\rightarrow$ archivos tocados $\rightarrow$ tests}.
\end{itemize}

\end{document}
